\chapter{Introduction}
\label{ch:introduction}

\section{Motivation}

Ant scans are valuable because they can turn specimens that are difficult to
measure repeatedly into a source of comparable three-dimensional traits. Once
anatomically corresponding models are available, quantities such as head
width, Weber's length, and appendage proportions can be extracted
systematically rather than measured by hand under a microscope. The practical
promise is therefore not only better mesh fitting, but a way to study
biodiversity at population scale.

In the first week of July 2026, ten preprocessed ant scans were passed
through the registration pipeline. All ten optimisation runs converged
cleanly: the loss curves showed no divergence, and the numerical objective
decreased exactly as expected. By the numerical signals available to the pipeline, every run had
succeeded.

Four of the resulting models were not recognisable as anatomically meaningful
ants. Legs had collapsed into the gaster; in some fits, the head had been
absorbed into the thorax. More troublingly, several of these anatomically
incorrect fits achieved a \emph{lower} loss than the runs that had visibly
succeeded.

The optimiser had succeeded. The registration had failed.

\findingmark{Established here; recurs throughout this report as \F{F1}.}%
This discrepancy is the starting point of the internship.\footnote{Recorded
as \F{F1} in the findings ledger (Section~\ref{sec:ledger}) and the direct
motivation for the correspondence instrument built later
(\S\ref{sec:synthetic-instrument}).} The objective minimised during fitting
is a symmetric point-to-point distance between two surfaces. It quantifies
whether the fitted mesh \emph{covers} the target, not whether vertex $i$ of
the template has landed on the anatomical point vertex $i$ is defined to
represent. These two properties can diverge, and when they do the failure
produces no exception, no divergent loss curve, and no signal in the log
that anything is wrong: the pipeline's own numerical output cannot tell a
biologically correct fit from a plausible but wrong one.

The project began from a fundamental question: how can a parametric
mesh be fitted to imperfect animal scans if the pipeline has no objective that
disambiguates good correspondence from a plausible but wrong deformation? The
internship addressed this question by moving from a mesh-quality-centred
hypothesis to a more precise investigation of initialisation, correspondence,
and constrained optimisation. A method is not considered reliable simply
because it reaches a numerically stable minimum, but because it reaches a
biologically meaningful one.

\section{Scope of the internship}

The internship brief listed six assignments: unify the preprocessing
pipeline, stabilise a Shape-Diameter-Function-based segmentation, integrate
neural initialisation, design a hybrid registration framework, benchmark the
result, and deliver documentation. Read together, this list implies a
specific hypothesis: that input mesh quality is the limiting factor, and
that cleaning it will improve the fits.
\S\ref{sec:turn} describes how an early, direct test of that hypothesis
produced a null result, and how the investigation was redirected on the
basis of that evidence toward initialisation, sampling, anatomical
structure, correspondence, and, ultimately, the optimisation objective
itself: the sequence this report follows.

\section{Parametric shape modelling for non-human subjects}

Parametric mesh registration deforms a rigged, segmented template mesh until
it matches a target scan, governed by a shape vector $\bm\beta$, joint
rotations $\bm\theta$, a global scale $s$, and residual per-vertex
deformation $\mathbf d$ (Eq.~\ref{eq:registration}). Because every
vertex retains a fixed anatomical identity across all fits, a correctly
registered corpus places otherwise unrelated scans into a shared coordinate
system in which measurement and population-level comparison become possible
: this property, not visual resemblance, is what makes the technique
useful, and Chapter~\ref{ch:pipeline} traces its lineage from \term{SMPL}
through \term{SMAL} to \term{SMILify} in full. An ant presents six limbs,
two antennae, and a pair of mandibles, all thin and articulated, several
represented by only a few dozen template vertices; scans from \term{scAnt}
photogrammetry~\cite{scant} or micro-CT commonly contain holes,
self-intersections, arbitrary orientation, and missing surface exactly at
the structures of greatest anatomical interest. Existing registration
pipelines, designed for clean, watertight meshes in near-canonical poses,
were not built to survive input of this kind: the problem this report
addresses.

\section{Project state at the outset}

\term{SMILify} is built on \term{SMALify}~\cite{smalify}, Biggs et al.'s
optimisation-based framework for fitting the SMAL quadruped model, itself
built on SMAL~\cite{smal}; its stated ambition is to turn any rigged mesh
into a SMAL-compatible parametric model regardless of skeletal topology, with
\term{SMIL} (Skinned Multi-Insect Linear) as the current insect-focused
instance. At the outset, a Shape Diameter Function module intended for
part-aware segmentation had not been validated on incomplete meshes, and a
\term{PointNet}-based initialisation predictor was implemented but
unevaluated; of roughly 380 raw scans in the available corpus, only 86
(around 23\%) were curated as clean enough for registration, indicating the
scale of the data-quality problem the brief was written against. A
longer-term ambition beyond this internship's scope, tracked internally as
\term{SMILify} Gen2, is building a shape space from the full scan collection
rather than a curated subset, something that can only be trusted once the
registrations feeding it are individually reliable, which is the
prerequisite this report's investigation addresses.


\section{Problem statement and research questions}
\label{sec:problem}

Registration recovers, for a scan $T$, the latent variables of a parametric
mesh $M$:
\begin{equation}
(\hat{\bm\beta},\hat{\bm\theta},\hat s,\hat{\mathbf d})
 =\arg\min_{\bm\beta,\bm\theta,s,\mathbf d}\;
 \mathcal{L}\big(M(\bm\beta,\bm\theta,s,\mathbf d),\,T\big),
\label{eq:registration}
\end{equation}
with shape $\bm\beta$, articulated pose $\bm\theta$, scale $s$ and residual
per-vertex deformation $\mathbf d$. The scan is only a set of surface
samples. It carries no vertex identities, no joint locations and no
labelling of anatomical structures, yet the aim is to recover exactly those
quantities. Equation~\eqref{eq:registration} is a non-convex problem whose
objective observes only surfaces (\S\ref{sec:registration}), so several
different latent configurations can yield nearly the same surface.

\begin{tcolorbox}[invstyle,title={Formal problem statement}]
The mapping from observed surface to anatomical state is not guaranteed to
be \emph{identifiable} from a surface-only objective alone: distinct
$(\bm\beta,\bm\theta,\mathbf d)$ can generate near-identical surfaces while
differing in every anatomical quantity that matters. Every mechanism this
report investigates (deformation absorbing pose error, local descriptors
that do not deploy densely, supervision that does not propagate) is a
different symptom of this one ambiguity.
\end{tcolorbox}

This is the central issue of the report, studied through four research
questions:

\begin{enumerate}[nosep,label=\textbf{RQ\arabic*.}]
\item Is input mesh quality the dominant limit on registration accuracy?
\item Does a good surface fit imply a correct anatomical fit?
\item Which component (initialisation, sampling, deformation freedom,
correspondence, or the objective) limits anatomical recovery, and where
along the body?
\item When a parameter is recovered, when is it a valid biological
measurement?
\end{enumerate}

\begin{figure}[!htb]
\centering
\begin{tikzpicture}[node distance=5mm and 5mm, every node/.style={font=\scriptsize},
 n/.style={draw,rounded corners=2pt,minimum height=7mm,align=center,fill=white,inner sep=2pt},
 arr/.style={-{Stealth[length=1.5mm]}}]
\node[n,fill=gray!10] (h0) {``The scans are the problem''};
\node[n,below=of h0] (pp) {Preprocessing test};
\node[n,below=of pp,fill=gray!10] (null) {Null result (\F{F2})};
\node[n,below=of null] (q) {Then what is failing?};
\node[n,below left=8mm and 30mm of q] (init) {Initialisation\\distal helps, coxa does not (\F{F4})};
\node[n,below=8mm of q] (corr) {Correspondence\\local signal exists, global deploy collapses};
\node[n,below right=8mm and 30mm of q] (obj) {Objective\\surface $\neq$ anatomy (\F{F1}); no joint term (\F{F6})};
\node[n,below=14mm of corr,fill=gray!10] (ct) {Counterfactual tests (oracles) + corrected expert ground truth (\F{F7})};
\node[n,below=of ct] (surv) {What actually survives?};
\node[n,below left=6mm and 22mm of surv] (s1) {Surface: reliable as coverage only};
\node[n,below=6mm of surv] (s2) {Anatomy: partial, structured recovery (\F{F8})};
\node[n,below right=6mm and 22mm of surv] (s3) {Measurement: trait-specific validity};
\node[n,below=10mm of s2,fill=gray!10] (open) {Open problem: a globally consistent anatomical correspondence};
\draw[arr](h0)--(pp); \draw[arr](pp)--(null); \draw[arr](null)--(q);
\draw[arr](q)--(init); \draw[arr](q)--(corr); \draw[arr](q)--(obj);
\draw[arr](init)|-(ct); \draw[arr](corr)--(ct); \draw[arr](obj)|-(ct);
\draw[arr](ct)--(surv); \draw[arr](surv)--(s1); \draw[arr](surv)--(s2); \draw[arr](surv)--(s3);
\draw[arr](s1)|-(open); \draw[arr](s2)--(open); \draw[arr](s3)|-(open);
\end{tikzpicture}
\caption{The investigation at a glance: from the initial mesh-quality
hypothesis, through its rejection, to a three-way diagnosis and what each
level of validity turned out to support.}
\label{fig:glance}
\end{figure}

\section{The redirection of the project}
\label{sec:turn}

The first weeks tested the hypothesis implicit in the brief (RQ1): a rebuilt
preprocessing pipeline was compared against the existing one with
registration held fixed. The two tied (\F{F2}, \S\ref{sec:prep-null}).
That null result located the problem elsewhere and redirected the internship
toward a question the pipeline had no instrument to answer: whether the
optimiser finds the \emph{right} correspondence, not merely a close surface.
Building that instrument from synthetic ground truth
(\S\ref{sec:synthetic-instrument}), and later an independent expert
benchmark (\S\ref{sec:skeleton-underdetermined}) after correcting a
coordinate-frame error in an early annotation set (\F{F7}), is the thread the
rest of the report follows.

\section{Established results}
\label{sec:ledger}

Results that recur are registered once here and cited by identifier
(\F{F1}, \F{F2}, and so on); evidence is in the chapter each points to. Table~\ref{tab:evidence}
(Chapter~\ref{ch:conclusion}) gives the evidence status of every headline
claim in the report.

\printfindingledger

\section{Contributions}

Existing SMILify provided the parametric representation and a baseline
optimisation-based fitting pipeline. This internship's contribution is the
diagnostic and validation machinery built around it, and what that machinery
established.

\textbf{Scientific contribution.} A diagnostic framework, evidenced
experimentally rather than asserted, showing that surface fit, anatomical
correspondence, parameter recovery and measurement validity are four
distinct properties that do not track one another: surface proximity does
not imply anatomical correspondence (\F{F1}); regional identity and
within-region localisation are distinct problems (\F{F3}); initialisation
sensitivity is anatomically heterogeneous (\F{F4}); the objective has no
explicit joint observation term and expert supervision does not propagate
spatially (\F{F6}); error is structured by joint and specimen rather than by
seed (\F{F8}); local learned correspondence does not imply dense global
correspondence (Chapter~\ref{ch:structure}); and accurate population-level
measurements can coexist with unreliable individual-level correspondence
(\S\ref{sec:scale-dependence}).

\textbf{Algorithmic contributions.} Chain-aware anatomical assignment; an
SDF-based anatomical field; a PointNet-style pose initialiser; localised
graduated non-convexity; a distal-protected sampling scheme; and learned
within-part correspondence descriptors, each evaluated against a
pre-registered criterion (Table~\ref{tab:evidence}), not assumed successful.

\textbf{Evaluation contributions.} A synthetic round-trip correspondence
benchmark; an adversarial metric test; a three-level evaluation protocol;
counterfactual (oracle) interventions for causal bottleneck localisation; a
corrected, frozen expert Joint Alignment Benchmark (11 specimens, 396 fits);
variance decomposition; and leave-one-region-out supervision.

\textbf{Systems and infrastructure contributions.} A configurable
preprocessing pipeline with Alpha Wrap and ManifoldPlus handling and
mesh-quality classification; integration of joint-limit priors into the 3D
registration path; a corrected staged fitting schedule; and reproducible
large-scale ablation studies across two HPC systems (JURECA, RWTH H100) with
fixed seeds, preserved checkpoints, configuration-driven experiments and
hash-frozen evaluation artifacts (\S\ref{sec:infrastructure}).

\textbf{What is open.} Within-part localisation on thin articulated
structures; the anterior deformation defect that the scale barrier did not
resolve; supplying joint information without expert annotation; and dense,
globally consistent learned correspondence. Claims not reproducible in the
final environment are marked provisional or excluded (Table~\ref{tab:evidence}).

\section{Structure of this report}

Chapter~\ref{ch:host} gives the host context and Chapter~\ref{ch:related}
positions the work against existing literature on parametric models,
point-cloud registration, shape correspondence, and defective mesh repair.
Chapter~\ref{ch:pipeline} defines the model, objective and pipeline
mathematically, and Chapter~\ref{ch:method} sets out the experimental
methodology. Chapter~\ref{ch:structure} reports the algorithmic
investigations and Chapter~\ref{ch:diagnostics} the mechanistic diagnosis of
the registration bottleneck. Chapter~\ref{ch:interventions} evaluates
candidate interventions against that diagnosis, Chapter~\ref{ch:validation}
is the move from registration to measurement, and Chapter~\ref{ch:conclusion}
the discussion, limitations and future work.
